\clearpage\thispagestyle{empty}
\pdfbookmark[0]{Sources for the October 1802 Ceres Report}{ceres180210-sources}
\begin{center}\textbf{Sources for the October 1802 Ceres Report}\end{center}
\noindent This AI translation follows Franz Xaver von Zach’s report on Ceres as excerpted in \href{https://gdz.sub.uni-goettingen.de/id/PPN236007467}{Gauss’s \textit{Werke}, vol.~VI} (Göttingen, 1874), pp.~224--225. Von Zach quotes Johann Friedrich Wurm and presents Carl Friedrich Gauss’s calculation of Jupiter’s perturbations of Ceres. Wurm’s italic Latin \textit{quid pro quo} refers to mistaking one thing for another.
\par\smallskip\noindent The \href{https://api.digitale-sammlungen.de/iiif/image/v2/bsb10538598_00392/full/1450,/0/default.jpg}{October 1802 first printing} is article XXXVIII of the \textit{Monatliche Correspondenz}, vol.~VI, pp.~382--389. The collected passage corresponds to the \href{https://api.digitale-sammlungen.de/iiif/image/v2/bsb10538598_00397/full/1450,/0/default.jpg}{prose and annual entries on p.~387} and the \href{https://api.digitale-sammlungen.de/iiif/image/v2/bsb10538598_00398/full/1450,/0/default.jpg}{periodic equations on p.~388}.
\par\smallskip\noindent The first printing calls \(0{,}000005909\) the annual increase in eccentricity. The collected edition calls \(0.000005909\) the annual decrease. In the journal, the argument of the longitude term \(+56{,}69\) ends with \(11^\circ\ 24'\ 1'\). The corresponding field in the \textit{Werke} is \(11\ 24\ 1''\). The journal gives \(23^\circ\ 42'\ 43''\) in the radius-vector term \(+0{,}0002469\), where the \textit{Werke} give \(23\ 42\ 53''\). Both translations follow the collected text.
\par\smallskip\noindent Translation and typesetting produced by AI. Journal scans from the Bayerische Staatsbibliothek.
\clearpage
